\documentclass[12pt]{article}
\usepackage[margin=1in]{geometry}
\usepackage{enumitem}
\usepackage{titlesec}

\titleformat{\section}{\large\bfseries}{}{0em}{}
\titlespacing*{\section}{0pt}{1em}{0.3em}

\setlength{\parindent}{0pt}
\setlength{\parskip}{0.5em}

\begin{document}

\begin{center}
    {\LARGE \textbf{Slice \& Dice}}\\[0.3em]
    {\large CIS 5680 --- Group Game Project 2}
\end{center}

\section{High Concept}

A VR blade-combat game where players wield an arsenal of edged weapons to carve through waves of oncoming targets with precision cuts, chaining rhythmic combos while dodging bombs and deflecting shields in an arena of physics-driven destruction.

\section{Genre}

First-person VR action-arcade with rhythm and precision-slicing mechanics.

\section{Features}

\begin{itemize}[leftmargin=2em]
    \item \textbf{Precision directional slicing.} Targets display highlighted cut lines at specific angles. Matching the angle exactly earns bonus points, while sloppy cuts still count but score less --- rewarding skill without punishing beginners.

    \item \textbf{Physics-based destruction.} Every slice produces physically simulated pieces that split, tumble, and shatter realistically. Cutting a melon in half feels different from cleaving through a wooden crate, giving each target a satisfying material identity.

    \item \textbf{Multiple weapon types.} Players choose from a katana, dual daggers, or a heavy cleaver, each with distinct swing speed, blade width, and feel. The katana rewards clean single strokes; daggers excel at rapid combos; the cleaver powers through tough targets.

    \item \textbf{Combo and rhythm system.} Consecutive accurate slices build a combo multiplier. A visual and audio beat guides the rhythm --- staying on tempo amplifies the multiplier, turning high-level play into a flowing, almost musical experience.

    \item \textbf{Hazard targets.} Not everything should be sliced. Bombs explode on contact, shields deflect the blade back, and flash grenades blind the player momentarily. Reading the incoming wave and reacting --- slice, dodge, or deflect --- keeps the player mentally engaged.

    \item \textbf{Dodge and deflect mechanics.} Players physically lean, duck, and sidestep to avoid hazards, and can use the flat of their blade to bat certain objects away. The full-body movement makes every round a workout.

    \item \textbf{Escalating waves.} Each wave increases in speed, density, and hazard mix. Early waves teach the mechanics; later waves demand split-second pattern recognition and weapon switching under pressure.

    \item \textbf{Arena variety.} Different arenas change how targets approach --- a dojo sends them from the front, a colosseum launches them from all directions, and a zero-gravity arena lets sliced pieces float, adding visual spectacle.
\end{itemize}

\section{Player Challenge / Motivation}

The player is a blade master facing an endless onslaught of targets. The motivation is score mastery: perfecting cut angles, maintaining combos, and surviving deeper into the waves. Leaderboards and personal bests drive replayability, while the physical satisfaction of a clean slice provides moment-to-moment reward.

\section{Design Goals}

\textbf{Visceral satisfaction.} Every swing should feel impactful. The combination of haptic feedback, physics-driven destruction, and audio cues makes slicing deeply satisfying on a sensory level.

\textbf{Easy to pick up, hard to master.} Swinging a blade is intuitive for any VR player. But chaining precision cuts at tempo while dodging hazards creates a skill ceiling that keeps experienced players coming back.

\textbf{Physical engagement.} The game is designed to get players moving --- reaching, swinging, ducking, and sidestepping. Sessions feel active and energizing rather than passive.

\section{Unique Selling Points}

\begin{itemize}[leftmargin=2em]
    \item \textbf{Angle-matched precision cutting} goes beyond simple ``hit the target'' mechanics. The directional requirement adds a layer of skill and intentionality to every swing that sets it apart from existing VR slicers.
    \item \textbf{Weapon variety with meaningful differences} gives players distinct playstyles to master rather than a single default blade, adding depth and replayability.
    \item \textbf{Rhythm-combo fusion} merges the satisfaction of a rhythm game with the physicality of melee combat, creating a genre hybrid that appeals to fans of both.
\end{itemize}

\end{document}
